% title: 讲义：高维代数簇 · 曲面的有理奇点
% renderer: btex
% legacy-title: 用户:Solution/ 讲义: 高维代数簇/曲面的有理奇点
% source: https://bananaspace.org/legacy/wiki/%E7%94%A8%E6%88%B7:Solution/_%E8%AE%B2%E4%B9%89:_%E9%AB%98%E7%BB%B4%E4%BB%A3%E6%95%B0%E7%B0%87/%E6%9B%B2%E9%9D%A2%E7%9A%84%E6%9C%89%E7%90%86%E5%A5%87%E7%82%B9
% summary: 从香蕉空间旧站迁移的原始 bTeX 页面。
\section{有理奇点}
\begin{definition}
    设$p \in S'$是孤立奇点，$S'$是正规曲面. 
    若存在一个奇点解消$\pi \colon S \to S'$，
    使得$R^{i}\pi_{*}\mathscr{O}_{S} = 0$对于任意$i > 0$成立，
    则称$p$是\textbf{有理奇点}，并称$\pi$是\textbf{有理解消}.
\end{definition}

\section{ADE曲线}
\begin{definition}
    设$C = \sum C_{i} \hookrightarrow S$是光滑曲面$S$上的
    例外曲线. 
    若诸$C_{i}^{2} = -2$且$C_{i} \cong \mathbb{P}^{1}$，
    则称$C$为ADE曲线.
\end{definition}